\tocsection{Centre of Mass of a Plane}
\question

Find the centre of mass of a plane.

\begin{parts}

\part Expected solution: Students use moments to find the centre of mass. They may only look at it in one dimension.

\begin{hints}

\begin{itemize}
    \item What is meant by the 'centre of mass' of an object?
    \item What forces are acting on the plane?
    \item Can you use moments?
\end{itemize}

\end{hints}

\begin{solution}

A good starting point is to recognise that the centre of mass is the point at which the total mass of the aircraft can be considered to act. 

If the aircraft is modelled as a collection of discrete components, and the mass and position of every component are known, then the centre of mass can be calculated as the mass-weighted average of their positions. 

For the longitudinal direction, 

\begin{equation}
x_{\mathrm{CM}} = \frac{\sum_i m_i x_i}{\sum_i m_i}
\end{equation}

where:
\begin{itemize} 
    \item \(m_i\) is the mass of component \(i\), 
    \item \(x_i\) is the position of component \(i\) relative to a chosen datum. 
\end{itemize} 

Similarly, the other two coordinates of the centre of mass are 

\begin{equation}
	y_{\mathrm{CM}} = \frac{\sum_i m_i y_i}{\sum_i m_i} \quad \text{and} \quad z_{\mathrm{CM}} = \frac{\sum_i m_i z_i}{\sum_i m_i}.
\end{equation}

Therefore, the three-dimensional centre of mass can be written as 

\begin{equation}
	\mathbf{r}_{\mathrm{CM}} = \frac{\sum_i m_i \mathbf{r}_i}{\sum_i m_i}.
\end{equation}

\end{solution}

\part Although the approach mentioned before is theoretically correct, calculating the centre of mass of a complete aircraft in this way would be difficult in practice. It would require knowledge of the mass and position of every significant component, including the following:

\begin{itemize} 
    \item fuselage and structural members, 
    \item wings and tail, \item engines, 
    \item landing gear, 
    \item fuel, 
    \item passengers and baggage, 
    \item electrical and mechanical systems. 
\end{itemize}

\end{parts}